% =====================================================================
%  ChromeFlashCardExtension — Context Handoff Document
%  Bien dich tren Overleaf: compiler mac dinh (pdfLaTeX) chay duoc ngay.
%  Neu tieng Viet hien sai, doi sang XeLaTeX (Menu > Compiler > XeLaTeX).
%  Preamble ben duoi tu dong tuong thich ca hai engine.
% =====================================================================
\documentclass[11pt,a4paper]{article}

% ---- Engine-aware font / encoding (pdfLaTeX hoac XeLaTeX/LuaLaTeX) ----
\usepackage{iftex}
\ifPDFTeX
  \usepackage[utf8]{inputenc}
  \usepackage[T5]{fontenc}   % T5 = bang ma tieng Viet (vntex)
  \usepackage{lmodern}
\else
  \usepackage{fontspec}      % XeLaTeX/LuaLaTeX: Unicode truc tiep
\fi

\usepackage[vietnamese]{babel}

% ---- Layout & typography ----
\usepackage[margin=2.4cm]{geometry}
\usepackage{xcolor}
\usepackage{underscore}      % cho phep go dau _ trong van ban thuong
\usepackage{booktabs}
\usepackage{tabularx}
\usepackage{array}
\usepackage{longtable}
\usepackage{enumitem}
\usepackage{listings}
\usepackage{titlesec}
\usepackage[most]{tcolorbox}
\usepackage{fancyhdr}
\usepackage{tikz}
\usetikzlibrary{arrows.meta, shapes.geometric}
\usepackage{hyperref}

% ---- Palette (dong bo voi mau accent cua chinh ung dung) ----
\definecolor{accent}{HTML}{2F4BBD}
\definecolor{accentsoft}{HTML}{EDF0FC}
\definecolor{ink}{HTML}{1A1D24}
\definecolor{muted}{HTML}{565F6D}
\definecolor{codebg}{HTML}{F4F5F7}
\definecolor{codeborder}{HTML}{E3E6EB}
\definecolor{ok}{HTML}{16794F}
\definecolor{warn}{HTML}{92600E}
\definecolor{bad}{HTML}{B4302C}

\hypersetup{
  colorlinks=true,
  linkcolor=accent,
  urlcolor=accent,
  citecolor=accent,
  pdftitle={ChromeFlashCardExtension - Context Handoff},
  pdfauthor={ChromeFlashCardExtension}
}

% ---- Section styling ----
\titleformat{\section}{\Large\bfseries\color{accent}}{\thesection.}{0.5em}{}
\titleformat{\subsection}{\large\bfseries\color{ink}}{\thesubsection}{0.5em}{}
\titlespacing*{\section}{0pt}{1.4em}{0.6em}

% ---- Header / footer ----
\pagestyle{fancy}
\fancyhf{}
\renewcommand{\headrulewidth}{0.4pt}
\renewcommand{\footrulewidth}{0pt}
\fancyhead[L]{\small\color{muted}ChromeFlashCardExtension}
\fancyhead[R]{\small\color{muted}Context Handoff}
\fancyfoot[C]{\small\color{muted}\thepage}

% ---- Code listing style ----
\lstdefinestyle{plain}{
  basicstyle=\ttfamily\footnotesize\color{ink},
  backgroundcolor=\color{codebg},
  frame=single,
  rulecolor=\color{codeborder},
  framesep=6pt,
  xleftmargin=6pt,
  xrightmargin=6pt,
  breaklines=true,
  breakatwhitespace=false,
  columns=fullflexible,
  keepspaces=true,
  showstringspaces=false,
  aboveskip=8pt,
  belowskip=4pt,
}
\lstset{style=plain}

% ---- Callout boxes ----
\newtcolorbox{note}[1][Ghi chú]{enhanced, breakable,
  colback=accentsoft, colframe=accent, coltitle=white,
  fonttitle=\bfseries\small, title=#1,
  boxrule=0.6pt, left=3mm, right=3mm, top=1.5mm, bottom=1.5mm}
\newtcolorbox{blocker}[1][Blocker]{enhanced, breakable,
  colback=bad!5, colframe=bad!70!black, coltitle=white,
  fonttitle=\bfseries\small, title=#1,
  boxrule=0.6pt, left=3mm, right=3mm, top=1.5mm, bottom=1.5mm}
\newtcolorbox{tip}[1][Mẹo]{enhanced, breakable,
  colback=ok!6, colframe=ok!70!black, coltitle=white,
  fonttitle=\bfseries\small, title=#1,
  boxrule=0.6pt, left=3mm, right=3mm, top=1.5mm, bottom=1.5mm}

% Small helper: inline code
\newcommand{\code}[1]{\texttt{#1}}

% Status tags
\newcommand{\tagok}{\textcolor{ok}{\textbf{Có}}}
\newcommand{\tagno}{\textcolor{bad}{\textbf{Chưa}}}
\newcommand{\tagpart}{\textcolor{warn}{\textbf{Một phần}}}

\setlist{itemsep=2pt, topsep=3pt, parsep=0pt}

% =====================================================================
\begin{document}

% ------------------------- TITLE PAGE --------------------------------
\begin{titlepage}
\centering
\vspace*{2.2cm}

{\color{accent}\rule{\linewidth}{2pt}}\\[0.5cm]
{\Huge\bfseries ChromeFlashCardExtension}\\[0.35cm]
{\Large\color{muted} Tiện ích học từ vựng + backend serverless trên AWS}\\[0.3cm]
{\color{accent}\rule{\linewidth}{2pt}}\\[1.4cm]

{\large\bfseries Tài liệu bàn giao ngữ cảnh (Context Handoff)}\\[0.4cm]
{\color{muted} Đọc tài liệu này là đủ để nắm bức tranh tổng thể của dự án:\\
sản phẩm, kiến trúc, API, dữ liệu, cấu hình, trạng thái và cách chạy thử.}\\[2.2cm]

\begin{tcolorbox}[width=0.8\linewidth, colback=accentsoft, colframe=accent,
  boxrule=0.8pt, arc=3pt]
\centering\small
\textbf{Loại dự án:} Đồ án / demo thực tập \quad\textbullet\quad
\textbf{Nền tảng:} Chrome Extension (MV3) + Node.js/Express + AWS\\[2pt]
\textbf{Nhánh tham chiếu:} \code{test-aws} \quad\textbullet\quad
\textbf{Snapshot:} 2026-07-14
\end{tcolorbox}

\vfill
{\color{muted}\small Sinh từ bộ tài liệu trong \code{docs/}. Xem mục cuối để biết đọc sâu ở đâu.}
\end{titlepage}

% ------------------------- TOC ---------------------------------------
\tableofcontents
\newpage

% =====================================================================
\section{Giới thiệu nhanh (đọc trong 60 giây)}

\begin{note}[TL;DR — Dự án này là gì?]
\textbf{ChromeFlashCardExtension} là một tiện ích Chrome giúp người học \emph{lưu từ vựng
ngay khi đang đọc web}, đồng bộ lên cloud, rồi \emph{ôn tập} và \emph{chơi game} trên web.
Backend là một ứng dụng Express duy nhất chạy được ở \textbf{hai chế độ}: local (file JSON)
và \textbf{AWS serverless} (Lambda + API Gateway + DynamoDB + S3 + Translate + CloudWatch).
Mục tiêu hiện tại: đưa backend lên AWS để làm demo thực tập dùng tối thiểu 3 dịch vụ AWS.
\end{note}

\noindent Sản phẩm gồm \textbf{4 bề mặt giao diện} dùng chung một backend:

\begin{enumerate}
  \item \textbf{Popup extension} — đăng nhập, thêm/sửa thẻ, đồng bộ, export.
  \item \textbf{Editor nội tuyến} (inline, chèn vào trang web qua context menu chuột phải).
  \item \textbf{Study Web} — học flashcard theo phiên, chấm mức nhớ (Again/Hard/Good/Easy), quản lý thư viện.
  \item \textbf{Game Web} — chơi solo đoán từ; có prototype realtime 1v1 nhưng \emph{chỉ chạy local}.
\end{enumerate}

% =====================================================================
\section{Sản phẩm dành cho ai và giải quyết điều gì}

\subsection{Người dùng chính}
\begin{itemize}
  \item \textbf{Người học}: đăng ký/đăng nhập; bôi đen từ trên trang web rồi mở editor bằng
        menu chuột phải; dịch sang tiếng Việt; sửa nghĩa/loại từ/nhóm rồi lưu; đồng bộ lên
        cloud; vào Study Web để ôn tập; vào Game Web để chơi; export JSON qua link tải tạm thời.
  \item \textbf{Người chấm / mentor}: xem kiến trúc dùng nhiều dịch vụ AWS có trách nhiệm rõ ràng;
        kiểm tra dữ liệu bền vững trên DynamoDB; xác nhận Lambda/API Gateway chạy độc lập với máy
        local; quan sát log/lỗi qua CloudWatch; đánh giá đánh đổi bảo mật/chi phí.
\end{itemize}

\subsection{Triết lý thiết kế}
\begin{itemize}
  \item \textbf{Offline-first}: thẻ luôn lưu được vào \code{chrome.storage.local} dù backend không chạy.
  \item \textbf{Một backend, hai chế độ}: cùng mã Express dùng cho cả local lẫn Lambda
        (qua \code{serverless-http}), tách biệt tầng lưu trữ (JSON so với DynamoDB).
  \item \textbf{Ranh giới tin cậy rõ ràng}: client không bao giờ được tin; backend luôn xác định
        chủ sở hữu dữ liệu từ JWT đã xác thực, không nhận \code{userId} từ client.
\end{itemize}

% =====================================================================
\section{Tính năng chính}

\begin{longtable}{@{}p{3.4cm} p{11.8cm}@{}}
\toprule
\textbf{Nhóm} & \textbf{Chi tiết} \\
\midrule
\endhead
Tài khoản & Đăng ký / đăng nhập bằng username + password; phiên đăng nhập giữ bằng JWT (hết hạn 7 ngày). \\
\addlinespace
Lưu thẻ khi đọc web & Bôi đen từ/cụm từ $\rightarrow$ chuột phải $\rightarrow$ ``Save as flashcard'' $\rightarrow$ editor hiện ngay trên trang (Shadow DOM, không đụng CSS trang gốc). \\
\addlinespace
Dịch tự động & Gọi \code{/api/translate}; trên AWS dùng Amazon Translate (mặc định dịch sang tiếng Việt), local dùng bản mock. \\
\addlinespace
Nhóm (category) & Tạo / liệt kê / xoá nhóm; xoá nhóm sẽ chuyển các thẻ về \code{Uncategorized}. \\
\addlinespace
Đồng bộ cloud & Upsert một chiều local $\rightarrow$ cloud theo \code{cardId}; không xoá thẻ trên cloud, không hai chiều. \\
\addlinespace
Study Web & Học theo phiên có xáo trộn, lật thẻ, chấm 4 mức nhớ (phím tắt 1--4), thanh tiến độ, tổng kết phiên; thư viện gom theo nhóm; CRUD thẻ. \\
\addlinespace
Game Web & Chơi solo (thử thách 10 thẻ hoặc chạy nước rút 30 giây), chấm điểm đúng/gần đúng/sai = 100/50/0. Tab realtime 1v1 là prototype local. \\
\addlinespace
Export & Backend ghi một file JSON vào bucket S3 riêng tư và trả về pre-signed URL (mặc định hết hạn sau 900 giây). \\
\addlinespace
Giao diện & Cả 4 bề mặt dùng chung một hệ design token, có chế độ sáng/tối (dark mode). \\
\bottomrule
\end{longtable}

% =====================================================================
\section{Kiến trúc tổng thể}

\subsection{Kiến trúc mục tiêu trên AWS (MVP)}

\begin{center}
\begin{tikzpicture}[
  font=\small,
  box/.style={draw=accent, thick, rounded corners, align=center,
              minimum height=9mm, minimum width=24mm, fill=accentsoft},
  cli/.style={draw=muted, thick, rounded corners, align=center,
              minimum height=9mm, minimum width=24mm, fill=gray!8},
  db/.style={draw=warn!80!black, thick, cylinder, shape border rotate=90,
             aspect=0.28, align=center, minimum height=12mm, minimum width=20mm,
             fill=warn!12},
  arr/.style={-{Latex[length=2.2mm]}, thick, draw=muted},
]
  % clients
  \node[cli] (ext) at (0,1.1) {Chrome\\Extension};
  \node[cli] (web) at (0,-1.1) {Study/Game\\Web (S3)};
  % core
  \node[box] (api) at (3.6,0) {API Gateway\\HTTP API};
  \node[box] (lam) at (7.2,0) {AWS Lambda\\Express};
  % data plane
  \node[db]  (ddb) at (11.0,0) {DynamoDB\\3 bảng};
  \node[box] (tr)  at (11.0,2.1) {Amazon\\Translate};
  \node[db]  (s3)  at (11.0,-2.1) {S3 Export\\(private)};
  \node[box] (cw)  at (7.2,-2.4) {CloudWatch};

  \draw[arr] (ext) -- (api);
  \draw[arr] (web) -- (api);
  \draw[arr] (api) -- (lam);
  \draw[arr] (lam) -- (ddb);
  \draw[arr] (lam) -- (tr);
  \draw[arr] (lam) -- (s3);
  \draw[arr] (lam) -- (cw);
  \draw[arr] (api) -- (cw);
\end{tikzpicture}
\end{center}

\noindent Ngoài ra, trình duyệt tải trực tiếp Study/Game (HTML/CSS/JS tĩnh) từ một
\textbf{bucket S3 website công khai} (chỉ chứa asset frontend). Khi bật \code{auto}
phát hiện ngôn ngữ, Amazon Translate còn kéo theo Amazon Comprehend
(\code{comprehend:DetectDominantLanguage}) như một dịch vụ hỗ trợ gián tiếp.

\subsection{Vì sao chọn serverless}
\begin{itemize}
  \item Không phải quản lý server/OS — hợp với demo nhỏ, scale-to-zero, tính tiền theo lượng dùng.
  \item Tái dùng nguyên ứng dụng Express hiện có qua \code{serverless-http}.
  \item DynamoDB khớp mẫu truy cập ``dữ liệu thuộc về từng user''.
  \item Mỗi dịch vụ có một trách nhiệm dễ trình bày trong báo cáo.
\end{itemize}

\subsection{Ranh giới tin cậy}
\begin{longtable}{@{}p{4.2cm} p{11cm}@{}}
\toprule
\textbf{Vùng} & \textbf{Thành phần} \\
\midrule
\endhead
Client (không tin) & extension, JS trên trình duyệt, \code{localStorage}, \code{chrome.storage}, mọi payload gửi lên. \\
Biên công khai & S3 website endpoint, API Gateway endpoint, pre-signed export URL. \\
AWS (tin cậy) & Lambda role, các bảng DynamoDB, lời gọi Translate/Comprehend, bucket export riêng tư. \\
\bottomrule
\end{longtable}

\begin{note}[Nguyên tắc vàng]
Không tin \code{userId}, vai trò, điểm số hay quyền sở hữu do client gửi lên. Backend luôn lấy
danh tính người dùng từ JWT đã xác thực (claim \code{sub}). Pre-signed URL là ``bearer capability'':
ai giữ URL trong thời gian còn hạn đều tải được — nên không log và không chia sẻ URL đó.
\end{note}

% =====================================================================
\section{Cấu trúc mã nguồn}

\begin{lstlisting}
ChromeFlashCardExtension/
|-- manifest.json            # Chrome Manifest V3 (permissions, content scripts)
|-- background.js            # Service worker: tao context menu, inject content script
|-- contentScript.js         # Editor noi tuyen tren trang web (Shadow DOM)
|-- popup.html/.css/.js       # Popup: auth, the local, sync, export
|-- extension-config.js       # Cau hinh API/Study URL cho extension
|
|-- backend/
|   |-- app.js               # Express app + routes (dung cho ca local va Lambda)
|   |-- server.js            # Khoi dong HTTP + WebSocket local (KHONG import trong Lambda)
|   |-- lambda.js            # Lambda handler (serverless-http)
|   |-- src/
|   |   |-- config.js               # Doc bien moi truong, chon che do
|   |   |-- localRepositories.js    # Luu tru JSON (chi development)
|   |   |-- dynamoRepositories.js   # Luu tru DynamoDB (che do cloud)
|   |   |-- translateService.js     # Mock hoac Amazon Translate
|   |   |-- exportService.js        # Export local hoac S3 pre-signed URL
|   |   |-- realtimeServer.js        # WebSocket local (in-memory, chi local)
|   |   |-- auth.js / validation.js / scoring.js / errors.js
|   |-- public/
|       |-- study/  (index.html, app.js, styles.css, config.js)
|       |-- game/   (index.html, app.js, styles.css, config.js)
|
|-- infra/template.yaml      # SAM starter (chua day du, con vai blocker)
|-- docs/                    # Bo tai lieu context AWS (11 file, xem muc cuoi)
`-- LOG.md                   # Nhat ky thay doi bat buoc cap nhat
\end{lstlisting}

\begin{note}[Nguồn sự thật (source of truth)]
Mã ở thư mục gốc và \code{backend/} là nguồn chính. Thư mục
\code{ChromeFlashCardExtension-test-aws-clean/} chỉ là bản sao thử nghiệm — \emph{đừng} coi là
nguồn chính và đừng đồng bộ song song. \code{infra/template.yaml} là starter, không phải bằng
chứng đã deploy AWS.
\end{note}

% =====================================================================
\section{Luồng người dùng}

\subsection{Hiện trạng (chạy local)}
\begin{lstlisting}
Boi den tren trang web
  -> Menu chuot phai (context menu)
  -> Editor duoc inject (contentScript.js)
  -> (tuy chon) POST /api/translate
  -> Luu vao chrome.storage.local
  -> Dang nhap o popup
  -> POST /api/sync  (origin = chrome-extension://<id>)
  -> Backend repository (JSON local)

Study Web: login/register -> flashcard & category API -> UI hoc/on tap
Game Web : REST API cho auth/cards -> WebSocket /realtime cho prototype 1v1
\end{lstlisting}

\subsection{Mục tiêu trên AWS}
Extension và Study/Game Web đều gọi HTTPS tới \emph{cùng} một API Gateway $\rightarrow$ Lambda
$\rightarrow$ DynamoDB, cộng thêm Amazon Translate và bucket export riêng tư. Trình duyệt tải
asset tĩnh trực tiếp từ bucket S3 website. Mọi dịch vụ đẩy log/metric về CloudWatch.

\begin{blocker}[Hai điểm bắt buộc sửa CODE trước khi deploy AWS]
\textbf{AUD-P0-07 — Translate trong editor:} \code{contentScript.js} hiện gọi
\code{/api/translate} trực tiếp. Trên trang web thật, request mang \emph{origin của trang đó}
(vd \code{https://en.wikipedia.org}), không phải origin extension, nên bị CORS chặn (403) khi
allowlist chặt trên AWS. Phải định tuyến qua \code{background.js} (service worker chạy ở origin
\code{chrome-extension://<id>}). \emph{Không} được nới CORS thành \code{*}.\\[4pt]
\textbf{AUD-P0-08 — Tắt realtime:} đặt \code{REALTIME\_URL=""} là \emph{không đủ} — code fallback
sang \code{ws://<origin>/realtime} và sẽ lỗi trên S3. Phải tắt UI realtime ngay trong code.
\end{blocker}

% =====================================================================
\section{Hợp đồng API (REST)}

\textbf{Quy ước:} JSON in/out; route cần bảo vệ dùng header \code{Authorization: Bearer <JWT>};
thành công thường có \code{ok: true}; lỗi có dạng \code{\{ ok: false, error: "..." \}}; giới hạn
body 1 MB. Base URL local: \code{http://localhost:3000}; base URL AWS:
\code{https://<api-id>.execute-api.<region>.amazonaws.com}.

\begin{longtable}{@{}p{4.6cm} p{1.2cm} p{4.4cm} p{4.4cm}@{}}
\toprule
\textbf{Method / Path} & \textbf{Auth} & \textbf{Request} & \textbf{Response} \\
\midrule
\endhead
\code{GET /api/health} & Không & --- & \code{\{ok, service\}} \\
\code{POST /api/auth/register} & Không & \code{\{username, password\}} & \code{\{ok, token, user\}} 201 \\
\code{POST /api/auth/login} & Không & \code{\{username, password\}} & \code{\{ok, token, user\}} \\
\code{GET /api/me} & Có & --- & \code{\{ok, user\}} \\
\code{GET /api/flashcards} & Có & --- & \code{\{ok, count, flashcards\}} \\
\code{POST /api/flashcards} & Có & flashcard & \code{\{ok, flashcard\}} 201 \\
\code{PUT /api/flashcards/:id} & Có & các field & \code{\{ok, flashcard\}} \\
\code{DELETE /api/flashcards/:id} & Có & --- & \code{\{ok, deletedId\}} \\
\code{GET /api/categories} & Có & --- & \code{\{ok, categories\}} \\
\code{POST /api/categories} & Có & \code{\{category\}} & \code{\{ok, category, categories\}} \\
\code{DELETE /api/categories/:cat} & Có & --- & category/count/list \\
\code{GET /api/study/random} & Có & \code{?category=} & \code{\{ok, flashcard, count\}} \\
\code{POST /api/sync} & Có & \code{\{flashcards:[...]\}} & count/created/updated/syncedAt \\
\code{POST /api/translate} & Có (AWS) & \code{\{word\}} hoặc \code{\{text,...\}} & bản dịch \\
\code{POST /api/export} & Có & \code{\{flashcards?\}} & \code{\{ok, fileName, downloadUrl\}} \\
\bottomrule
\end{longtable}

\noindent Ví dụ một số payload (đã bỏ dấu tiếng Việt trong khối mã cho tương thích trình biên dịch):

\begin{lstlisting}
// POST /api/auth/login  ->  response
{ "ok": true,
  "token": "<jwt>",
  "user": { "id": "<uuid>", "userId": "<uuid>",
            "username": "demo_user", "role": "user" } }

// POST /api/sync  ->  response
{ "ok": true, "count": 1, "created": 1, "updated": 0,
  "syncedAt": "2026-07-14T00:00:00.000Z" }

// POST /api/translate  ->  response (che do AWS)
{ "ok": true, "word": "resilient", "meaning": "kien cuong",
  "wordform": "unknown", "provider": "amazon-translate" }
\end{lstlisting}

\begin{note}[Ngữ nghĩa Sync cần nhớ]
\code{/api/sync} là \textbf{bulk upsert một chiều} local $\rightarrow$ cloud theo \code{cardId}.
Không phải full-replace, không lan truyền thao tác xoá ở local, và \emph{không atomic} — lỗi
giữa chừng có thể ghi được một phần (partial write). Retry cùng \code{cardId} thường cập nhật item
cũ nhưng timestamp/count có thể đổi.
\end{note}

% =====================================================================
\section{Mô hình dữ liệu}

\subsection{DynamoDB (chế độ cloud)}
\begin{longtable}{@{}p{3.1cm} p{4.6cm} p{7.3cm}@{}}
\toprule
\textbf{Bảng} & \textbf{Khoá} & \textbf{Ghi chú} \\
\midrule
\endhead
Users & PK: \code{username} & \code{userId} (UUID, = JWT \code{sub}), \code{passwordHash} (bcrypt), \code{role}. Không GSI. \\
\addlinespace
Flashcards & PK: \code{userId}, SK: \code{cardId} & Các field thẻ + \code{userId}. Lọc theo category hiện làm trong Lambda \emph{sau} khi Query. \\
\addlinespace
Categories & PK: \code{userId}, SK: \code{categoryName} & \code{createdAt}, \code{updatedAt}. \\
\bottomrule
\end{longtable}

\noindent Mẫu truy cập: Query toàn bộ thẻ/nhóm của user hiện tại; get/update/delete theo khoá kép.
Xác thực token: nạp lại user theo \code{username} rồi kiểm tra \code{userId === sub}.

\subsection{Hình dạng logic của một thẻ}
\begin{lstlisting}
{
  "id": "<= cardId, de tuong thich client>",
  "cardId": "<uuid hoac id san co>",
  "word": "resilient",
  "meaning": "kien cuong",
  "wordform": "adjective",
  "category": "IELTS",
  "sourceUrl": "https://example.test/article",
  "sourceTitle": "Demo article",
  "createdAt": "2026-07-14T00:00:00.000Z",
  "updatedAt": "2026-07-14T00:00:00.000Z",
  "syncedAt":  "2026-07-14T00:00:00.000Z"
}
\end{lstlisting}

\subsection{Lưu trữ phía trình duyệt}
\begin{longtable}{@{}p{3.2cm} p{5.0cm} p{6.8cm}@{}}
\toprule
\textbf{Nơi lưu} & \textbf{Key} & \textbf{Nội dung} \\
\midrule
\endhead
Extension & \code{flashcards} & Thẻ local (offline) \\
          & \code{flashcardCategories} & Danh sách nhóm \\
          & \code{flashcardAuth} & Token + info user \\
          & \code{flashcardTheme} & Lựa chọn sáng/tối của popup \\
\addlinespace
Study/Game & \code{flashcardStudyAuth} & Token + user (dùng chung nếu cùng origin) \\
           & \code{flashcardStudyTheme} & Lựa chọn sáng/tối (Study + Game dùng chung key) \\
\bottomrule
\end{longtable}

% =====================================================================
\section{Cấu hình và biến môi trường}

Chuyển từ local sang AWS chủ yếu là đổi biến môi trường. Bảng dưới là các biến quan trọng
(rút gọn từ \code{05\_API\_DATA\_CONFIG\_CONTRACTS.md}).

\begin{longtable}{@{}p{4.0cm} p{3.4cm} p{7.0cm}@{}}
\toprule
\textbf{Biến} & \textbf{Local} & \textbf{AWS} \\
\midrule
\endhead
\code{DATA_STORE} & \code{local} & \code{dynamodb} (đồng thời bật Translate thật, yêu cầu auth translate, tắt allow-all-origins) \\
\code{JWT_SECRET} & fallback yếu & \textbf{Bắt buộc} giá trị mạnh \\
\code{USERS_TABLE} & \code{FlashcardUsers} & Bắt buộc \\
\code{FLASHCARDS_TABLE} & \code{FlashcardCards} & Bắt buộc \\
\code{CATEGORIES_TABLE} & \code{FlashcardCategories} & Bắt buộc \\
\code{EXPORT_BUCKET} & rỗng & \textbf{Bắt buộc} (bucket riêng tư) \\
\code{ALLOWED_ORIGINS} & danh sách localhost & \textbf{Bắt buộc} — liệt kê chính xác, không \code{*} \\
\code{CORS_ALLOW_ALL} & có điều kiện & \textbf{Phải vắng / false} trên AWS \\
\code{SERVE_STUDY_STATIC} & \code{true} & \code{false} (S3 phục vụ frontend) \\
\code{AWS_REGION} & shell/profile & Runtime tự cấp — \textbf{không tự set} (reserved key) \\
\code{TRANSLATE_MAX_LENGTH} & \code{120} & \code{120} (chặn chi phí) \\
\bottomrule
\end{longtable}

\begin{note}[Vì sao ``chạy được ở local'' không chứng minh CORS đúng]
Trong \code{backend/src/config.js}, \code{allowAllOrigins} tự bật khi \code{DATA_STORE=local} và
chưa set \code{ALLOWED_ORIGINS}. Đó là lý do local không bao giờ lộ lỗi CORS. Đừng lấy ``local
chạy ngon'' làm bằng chứng cấu hình CORS chính xác cho AWS.
\end{note}

% =====================================================================
\section{Trạng thái hiện tại và các blocker}

\begin{longtable}{@{}p{6.6cm} p{2.4cm} p{5.4cm}@{}}
\toprule
\textbf{Khu vực} & \textbf{Trạng thái} & \textbf{Ghi chú} \\
\midrule
\endhead
Express chạy local & \tagok & \code{server.js}, repository JSON \\
Lambda handler & \tagok & \code{lambda.js}, \code{serverless-http} \\
DynamoDB repository & \tagok & 3 bảng, query theo \code{userId} \\
Amazon Translate & \tagpart & Có code; IAM còn thiếu quyền Comprehend cho \code{auto} \\
Translate trong editor & \tagno & Sẽ hỏng vì CORS trên AWS (AUD-P0-07) \\
Export S3 riêng tư & \tagok & Pre-signed URL 15 phút \\
API Gateway / SAM & \tagpart & Runtime \code{nodejs20.x} đã deprecated; template còn set reserved \code{AWS_REGION} \\
S3 Study/Game web & \tagpart & Có asset tĩnh; hạ tầng bucket chưa nằm trong SAM \\
Realtime multiplayer & \tagpart & Prototype WebSocket in-memory, không hợp Lambda HTTP API \\
Tắt realtime cho AWS & \tagno & \code{REALTIME_URL=""} không tắt được (AUD-P0-08) \\
Test tự động & \tagno & Mới có syntax check cho 3 file entry \\
AWS deploy đã xác nhận & \tagno & Chưa có bằng chứng trong repo \\
\bottomrule
\end{longtable}

\begin{blocker}[Không deploy trước khi xử lý các mục P0]
Trong 8 blocker P0: 6 mục chỉ cần đổi cấu hình hoặc chốt phạm vi (nâng runtime lên
\code{nodejs24.x}, thêm quyền IAM cho Translate, bỏ reserved \code{AWS_REGION}, gói artifact sạch,
chốt cấu hình URL, gắn nhãn realtime là ngoài phạm vi). Riêng \textbf{AUD-P0-07} và
\textbf{AUD-P0-08} \emph{bắt buộc sửa code} và không thể đóng bằng runbook.
\end{blocker}

% =====================================================================
\section{Chạy thử ở local (Getting Started)}

\textbf{Yêu cầu:} Node.js $\ge 18$ (máy dev đang dùng v24) và trình duyệt nhân Chromium
(Chrome/Brave/Edge).

\subsection{1) Khởi động backend}
\begin{lstlisting}
cd backend
npm install        # lan dau
npm run dev        # = node --watch server.js  ->  http://localhost:3000
\end{lstlisting}

\noindent Sau khi server chạy:
\begin{itemize}
  \item Study Web: \code{http://localhost:3000/study/}
  \item Game Web: \code{http://localhost:3000/game/}
  \item Tài khoản mẫu (seed local): \code{student} / \code{password123}
\end{itemize}

\begin{tip}[Sửa CSS/HTML của Study/Game]
\code{--watch} chỉ theo dõi file backend. Khi đổi asset trong \code{backend/public/}, không cần
restart server — chỉ cần hard-reload trình duyệt (Ctrl+Shift+R) để bỏ cache CSS cũ.
\end{tip}

\subsection{2) Nạp extension vào trình duyệt}
\begin{enumerate}
  \item Mở \code{chrome://extensions} (hoặc \code{brave://extensions}).
  \item Bật \textbf{Developer mode}.
  \item \textbf{Load unpacked} $\rightarrow$ chọn \emph{thư mục gốc} của dự án (nơi có \code{manifest.json}).
  \item Mỗi lần đổi code extension: bấm nút \textbf{Reload} trên card của extension,
        rồi \textbf{F5 lại tab web} đang mở để nạp content script mới.
\end{enumerate}

\begin{note}[Content script không chạy trên trang nội bộ trình duyệt]
Menu chuột phải ``Save as flashcard'' luôn hiện (do service worker tạo), nhưng editor chỉ hiện
được trên trang \code{http://} / \code{https://} \emph{thật}. Trên trang New Tab, \code{chrome://},
\code{brave://}, Web Store... trình duyệt chặn cứng mọi extension. Muốn thử: mở một trang thật
(vd Wikipedia), bôi đen một từ rồi chuột phải.
\end{note}

% =====================================================================
\section{Bảo mật, vận hành và chi phí (tóm tắt)}

\begin{itemize}
  \item \textbf{Bí mật:} không commit AWS credentials, JWT secret production hay \code{.env}.
        Lambda nhận quyền qua execution role, không nhúng access key.
  \item \textbf{Least privilege:} Lambda role chỉ đủ quyền cho 3 bảng, prefix object export,
        Translate (và Comprehend nếu dùng \code{auto}).
  \item \textbf{Export bucket:} bật Block Public Access; URL thô phải trả 403, chỉ pre-signed URL tải được.
  \item \textbf{CORS:} không dùng \code{*} ở bản demo cuối; allowlist chính xác theo origin.
  \item \textbf{Log an toàn:} không log JWT, password, full pre-signed URL hay nội dung nhạy cảm;
        đặt retention hữu hạn (khuyến nghị 14 ngày cho demo).
  \item \textbf{Alarm tối thiểu:} Lambda Errors và API Gateway 5xx; thêm Throttles nếu có thời gian.
  \item \textbf{Chi phí:} đặt budget nhỏ + billing alert; DynamoDB khởi điểm 1 RCU/1 WCU cho demo,
        không phải SLA production.
\end{itemize}

\begin{note}[Quyền riêng tư]
Thẻ có thể chứa \code{sourceUrl}/\code{sourceTitle} — tức lịch sử đọc web của người dùng. Vì vậy
demo phải dùng dữ liệu giả; trước khi lên production cần có privacy notice và chính sách
lưu trữ/export/xoá.
\end{note}

% =====================================================================
\section{Đọc sâu ở đâu (bản đồ tài liệu)}

Toàn bộ chi tiết nằm trong thư mục \code{docs/} (tiếng Việt). Thứ tự đọc khuyến nghị:

\begin{longtable}{@{}p{6.4cm} p{8.4cm}@{}}
\toprule
\textbf{File} & \textbf{Nội dung} \\
\midrule
\endhead
\code{01_PROJECT_CONTEXT.md} & Sản phẩm, người dùng, phạm vi, cấu trúc repo \\
\code{02_REQUIREMENTS.md} & Yêu cầu chức năng/phi chức năng, tiêu chí chấp nhận \\
\code{03_CURRENT_STATE_AUDIT.md} & Hiện có gì, bằng chứng, khoảng trống \\
\code{04_AWS_ARCHITECTURE.md} & Kiến trúc mục tiêu, quyết định (ADR), ranh giới MVP \\
\code{05_API_DATA_CONFIG_CONTRACTS.md} & API, DynamoDB, storage, biến môi trường \\
\code{06_MIGRATION_PLAN.md} & Thứ tự thay đổi an toàn và các cổng kiểm soát \\
\code{07_MANUAL_DEPLOYMENT_RUNBOOK.md} & Triển khai bằng AWS Console \\
\code{08_SECURITY_OPERATIONS_COST.md} & Bảo mật, log, alarm, backup, chi phí \\
\code{09_TEST_AND_ACCEPTANCE.md} & Ma trận test và bằng chứng cần lưu \\
\code{10_IMPROVEMENT_BACKLOG.md} & Backlog ưu tiên, definition of done \\
\code{11_AGENT_HANDOFF.md} & Quy tắc cho người/agent tiếp theo \\
\bottomrule
\end{longtable}

\vspace{0.6cm}
\begin{center}
\color{muted}\small
--- Hết tài liệu bàn giao ngữ cảnh ---\\
Cập nhật cùng nhịp với \code{LOG.md} mỗi khi kiến trúc hoặc hợp đồng thay đổi.
\end{center}

\end{document}
